\section{Data Privacy \& GDPR}
\label{sec:ethics-legal-societal}
\label{sec:data-privacy-gdpr}
MERIT handles personal hiring data. The sections below summarise UK GDPR obligations as well as the ethical limitations of
the prototype along with its wider effects on technical hiring. We build upon the safeguards in Section~\ref{sec:bias-mitigation-safeguards}
and the evaluation in Chapter~\ref{ch:evaluation}.

If MERIT is to be deployed and used for real hiring in the UK, it would fall under data-protection law (UK GDPR) \cite{uk_dpa_2018}.
MERIT stores sensitive PII such as names, contact details and CV content (which in itself contains PII).
Alongside this, it stores LinkedIn content and GitHub activity, which could be considered personal data.
Alongside this, we also store generated outputs such as the match score, Shapley breakdowns and fraud flags.
In production, the employer using this system would be legally responsible for this data, rather than the author of MERIT themselves.
Hiring teams usually rely on contract law, or some form of business need to process applications. Here, candidates must be asked for
consent to process their data \cite{ico_employment_2023}.

MERIT also stores and displays data such as profile pictures and university names, which may be considered personal data. While this
may be masked in Blind Mode, it is still stored in the backend database. This data is likely to hint at background characteristics,
which employers must not discriminate on. To encourage appropriate behaviour, Blind Mode is enabled by default, and it is the recruiter's
own choice if they wish to disable it. If MERIT is to be used in the hiring process, employers should publish a clear notice of what 
data is being collected, how it will be used and encrypted. There will also need to be a deletion date, and agreements with any third-party 
hosts. There also should be options for candidates to opt-out of the data collection and use, and any data gathered for this process
should not be reused elsewhere, without further consent.

Article~22 of the GDPR protects individuals from decisions made entirely by algorithms without human intervention, giving candidates
the right to not be subject to a decision based solely on automated processing if that decision produces legal effects or significantly 
impacts them \cite{eu_gdpr_2016}. The outcome of a candidate screen is most certainly a significant effect on a candidate's future, and thus, we should 
never rely on MERIT as an autonomous screening tool. MERIT is designed to be a decision-support tool first and foremost, and a recruiter
is encouraged to review outcomes, which is why we have the pillar of explainability. We provide recruiters the tools to challenge
our algorithmically-generated scores. Public GitHub and LinkedIn pages still fall under GDPR as personal data as well, since
these profiles store full names, email addresses, phone numbers and other such contact details.

During FYP development, we collected data from at most 30 candidate profiles (CV, GitHub, and LinkedIn) from contacts who gave informed consent,
including permission to use their data as a basis for synthetic evaluation personas. For this research we used the Apify actor
\texttt{harvestapi/linkedin-profile-scraper}; evaluation data were otherwise synthetic or anonymised, stored in Supabase with API keys in local
environment variables only. Consent makes this research more reasonable, but it does not make Apify-style collection acceptable for production:
deployments should use licensed LinkedIn APIs, privacy impact assessments, and subject access and deletion processes \cite{ico_employment_2023}.
Legal compliance still depends on how the employer operates the tool.

\subsection{UK employment law and AI governance}
\label{sec:uk-ai-governance}

If employers are to use MERIT within the UK, they must be aware of the Equality Act~2010~\cite{uk_equality_act_2010}.
There is a prohibition on indirect discrimination on protected characteristics such as gender, race and age. MERIT's scoring engine
does not discriminate on these characteristics, and Blind Mode is enabled by default to mask this information, however, it is the
employer's choice if they wish to disable Blind Mode. While we do have a study for measuring prestige-proxy bias, it does not
constitute the equality impact assessment that would be required for such a tool before we could deploy it to production.

The ICO's employment guidance~\cite{ico_employment_2023} expects both transparency as well as human oversight when it comes to the 
use of AI in hiring. The UK does not have a specific overarching AI act, but rather has sector-led approaches to this regulation,
rather than an omnibus statute (like the EU AI Act \cite{eu_ai_act_2024}). While MERIT's glass-box audit trail defaults to Blind
Mode to protect PII, the prototype still lacks representative dataset governance and the previously mentioned impact assessment.
\footnote{Note that if MERIT were used to screen candidates who are
based in the EU, Regulation (EU) 2024/1689 (the EU Artificial Intelligence Act) would also apply, and Annex~III classifies recruitment
screening as high-risk~\cite{eu_ai_act_2024}. This report focuses only on UK deployment.}
